\chapter{Quantum States, Operators, and Hamiltonians}
\index{quantum state|textbf}
\index{operator!definition|textbf}
\index{Hamiltonian!definition|textbf}

\label{ch:training-quantum-operators}

\chapterstatus{pedagogical quantum-mechanical foundation. The finite example is
illustrative and establishes no material Hamiltonian, calculated scientific
result, validation, or uncertainty quantification.}

\section{Learning goals}

After working through this chapter, a reader should be able to

\begin{enumerate}
  \item distinguish an abstract quantum state from its coordinates;
  \item state the domain-qualified definition of a linear operator;
  \item interpret expectation values and Hamiltonian eigenproblems;
  \item distinguish an operator identity from a matrix identity;
  \item solve the one-dimensional particle-in-a-box eigenproblem from its
        boundary conditions;
  \item apply Born's interpretation to position and energy measurements; and
  \item translate that continuum problem into a finite representation without
        identifying the finite matrix with the abstract operator itself.
\end{enumerate}

\section{State spaces and Dirac notation}
\index{Dirac notation|textbf}
\index{Hilbert space}
\index{expectation value}

A one-particle quantum model begins by declaring a Hilbert space \(\mathcal H\).
A quantum state is represented by a ket
\(\lvert\psi\rangle\in\mathcal H\). Its corresponding bra
\(\langle\psi\rvert\) is the linear functional associated with the ket by the
Hilbert-space inner product. For
\(\lvert\phi\rangle,\lvert\psi\rangle\in\mathcal H\),
\begin{equation}
  \langle\phi\vert\psi\rangle,
  \qquad
  \lVert\psi\rVert
  =\sqrt{\langle\psi\vert\psi\rangle}.
\end{equation}
A normalized state satisfies
\(\langle\psi\vert\psi\rangle=1\).

If \(\{\lvert\phi_a\rangle\}\) is an orthonormal basis, then
\begin{equation}
  \lvert\psi\rangle
  =\sum_a c_a\lvert\phi_a\rangle,
  \qquad
  c_a=\langle\phi_a\vert\psi\rangle.
\end{equation}
The ket is the abstract state; the coefficients \(c_a\) are its coordinates in
the selected basis. Changing the basis changes the coordinates without, by
itself, changing the state. This is the quantum-mechanical use of the
finite-representation distinction developed in
Chapter~\ref{ch:training-finite-representations}.

\section{Linear operators and domains}
\index{domain!operator|textbf}

\begin{definition}[Linear operator]
Let \(\mathcal H\) be a Hilbert space of quantum states. A linear operator
\(\hat A\) is a map
\begin{equation}
  \hat A:\mathcal D(\hat A)\subseteq\mathcal H\longrightarrow\mathcal H,
\end{equation}
with a domain \(\mathcal D(\hat A)\) on which its action is defined.
\end{definition}

The domain is part of the operator. Writing
\(\hat A\lvert\psi\rangle\) is meaningful only when
\(\lvert\psi\rangle\in\mathcal D(\hat A)\). Differential expressions that look
identical can define different operators when supplied with different domains
or boundary conditions.

For a normalized state in the domain of a self-adjoint observable \(\hat A\),
the expectation value is
\begin{equation}
  \langle\hat A\rangle_{\psi}
  =\langle\psi\rvert\hat A\lvert\psi\rangle.
\end{equation}
Self-adjointness is an operator property, not merely a visual symmetry of one
array. In a finite orthonormal representation it implies a Hermitian matrix, but
that finite test does not reconstruct unrecorded domain or boundary
information.

\section{Hamiltonians and additive decompositions}
\index{Hamiltonian!additive decomposition}

\begin{definition}[Hamiltonian]
A Hamiltonian \(\hat H\) is the operator that generates the dynamics of the
selected physical model. For the elementary nonrelativistic decomposition used
here,
\begin{equation}
  \hat H=\hat T+\hat V,
\end{equation}
where \(\hat T\) is the kinetic-energy operator and \(\hat V\) is the potential
or interaction operator.
\end{definition}

For every state in the common domain,
\begin{equation}
  \hat H\lvert\psi\rangle
  =\bigl(\hat T+\hat V\bigr)\lvert\psi\rangle.
\end{equation}
This is an equality of operator actions on states. It is not yet a finite matrix
equality, because no finite space or basis has been selected.

The stationary eigenproblem is
\begin{equation}
  \hat H\lvert\psi_n\rangle
  =E_n\lvert\psi_n\rangle.
\end{equation}
Energy eigenvalues are important outputs, but they do not alone identify the
state correspondence or the action of \(\hat H\) on arbitrary superpositions.
A list of energies is therefore not a complete Hamiltonian representation.

\section{Position-space representation}
\index{position-space representation|textbf}
\index{coordinate representation|textbf}

Let \(\langle\mathbf r\rvert\) be the generalized position bra and define the
wavefunction by
\(\psi(\mathbf r)=\langle\mathbf r\vert\psi\rangle\). The position-space
representation of the abstract operator equation is
\begin{equation}
  \langle\mathbf r\rvert\hat H\lvert\psi\rangle
  =\langle\mathbf r\rvert\hat T\lvert\psi\rangle
  +\langle\mathbf r\rvert\hat V\lvert\psi\rangle.
\end{equation}
For the nonrelativistic kinetic operator and a local scalar potential, this
becomes
\begin{equation}
  \langle\mathbf r\rvert\hat H\lvert\psi\rangle
  =-\frac{\hbar^{2}}{2m}\nabla^{2}\psi(\mathbf r)
  +V(\mathbf r)\psi(\mathbf r).
\end{equation}

The generalized coordinate bra \(\langle\mathbf r\rvert\) is not a finite-rank
projector. Coordinate representation, subspace projection, basis selection,
and model reduction are distinct operations.

A local multiplication operator in one representation need not remain diagonal
or local after projection or basis transformation. ``Kinetic'' and
``potential'' describe roles in the selected model; they do not authorize
subtraction of matrices assembled in incompatible spaces.

\section{The one-dimensional particle in a box}
\index{particle in a box!one-dimensional|textbf}
\index{infinite square well}
\index{Dirichlet boundary condition}

The one-dimensional particle in a box turns the preceding definitions into one
complete eigenproblem. Consider a particle of mass \(m\) confined between two
endpoints \(a<b\). The interior and its length are
\begin{equation}
  \Omega=(a,b),
  \qquad
  L=b-a>0.
\end{equation}
Its state space is \(L^2(a,b)\), the space of square-integrable complex
wavefunctions on the interval. The inner product is
\begin{equation}
  \langle\phi\vert\psi\rangle
  =\int_a^b\phi(x)^*\psi(x)\,\mathrm{d}x.
\end{equation}
A normalized state satisfies
\begin{equation}
  \int_a^b |\psi(x)|^2\,\mathrm{d}x=1.
\end{equation}

Inside the box, the potential is zero. Confinement is imposed at the closed
interval's endpoints through the Dirichlet boundary conditions
\begin{equation}
  \psi(a)=0,
  \qquad
  \psi(b)=0.
\end{equation}
The Hamiltonian is therefore the kinetic operator
\begin{equation}
  \hat H_{\mathrm{box}}
  =-\frac{\hbar^2}{2m}\frac{\mathrm{d}^2}{\mathrm{d}x^2}
\end{equation}
on a domain of sufficiently regular wavefunctions that satisfy those boundary
conditions. The boundary conditions are part of the operator. Reusing the same
second-derivative expression with periodic or Neumann boundary conditions would
define a different Hamiltonian.

The phrase ``infinite square well'' is common because one may imagine an
exterior potential that becomes infinitely large. In the operator used here,
there is no finite interior wall matrix to add. The confinement has already
entered through the interval and the domain. Appendix~\ref{app:particle-box}
develops this distinction formally and uses it for a separate retained-space
residual diagnostic~\cite{bonneau2001,angelone2024}.

\subsection{Solving the stationary equation by a characteristic equation}
\index{characteristic equation!differential equation}

The stationary Schrödinger equation is
\begin{equation}
  -\frac{\hbar^2}{2m}\frac{\mathrm{d}^2\psi}{\mathrm{d}x^2}=E\psi.
\end{equation}
For \(E>0\), define
\begin{equation}
  k^2=\frac{2mE}{\hbar^2}.
\end{equation}
The differential equation becomes
\begin{equation}
  \frac{\mathrm{d}^2\psi}{\mathrm{d}x^2}+k^2\psi=0.
\end{equation}

Use the exponential trial function
\begin{equation}
  \psi(x)=e^{rx},
\end{equation}
where \(r\) is initially an unknown number. Differentiation gives
\begin{equation}
  \frac{\mathrm{d}\psi}{\mathrm{d}x}=re^{rx},
  \qquad
  \frac{\mathrm{d}^2\psi}{\mathrm{d}x^2}=r^2e^{rx}.
\end{equation}
Substitution into the differential equation yields
\begin{equation}
  r^2e^{rx}+k^2e^{rx}=0.
\end{equation}
The exponential never vanishes, so divide by \(e^{rx}\):
\begin{equation}
  r^2+k^2=0.
\end{equation}
This polynomial is the characteristic equation of the constant-coefficient
ordinary differential equation. Its roots are
\begin{equation}
  r_{+}=ik,
  \qquad
  r_{-}=-ik.
\end{equation}
Therefore the complex exponential form of the general solution is
\begin{equation}
  \psi(x)=C_{+}e^{ikx}+C_{-}e^{-ikx}.
\end{equation}
Euler's identities give the equivalent real-function form
\begin{equation}
  \psi(x)=A\sin(kx)+B\cos(kx).
\end{equation}

For the interval \([a,b]\), it is convenient to measure distance from the left
endpoint. The same two-dimensional solution space can be written as
\begin{equation}
  \psi(x)
  =A\sin\!\bigl(k(x-a)\bigr)
   +B\cos\!\bigl(k(x-a)\bigr).
\end{equation}
Apply the left boundary condition:
\begin{equation}
  \psi(a)=B=0.
\end{equation}
The surviving solution is
\begin{equation}
  \psi(x)=A\sin\!\bigl(k(x-a)\bigr).
\end{equation}
At the right endpoint,
\begin{equation}
  \psi(b)=A\sin\!\bigl(k(b-a)\bigr)=A\sin(kL)=0.
\end{equation}
A nonzero eigenfunction has \(A\neq0\), so
\begin{equation}
  \sin(kL)=0,
  \qquad
  kL=n\pi,
  \qquad
  n=1,2,3,\ldots.
\end{equation}
Thus the permitted wave numbers are
\begin{equation}
  k_n=\frac{n\pi}{L}.
\end{equation}
Substitution into \(E=\hbar^2k^2/(2m)\) gives
\begin{equation}
  E_n
  =\frac{\hbar^2\pi^2n^2}{2mL^2},
  \qquad
  n=1,2,3,\ldots.
\end{equation}
There is no \(n=0\) state: it would produce the zero function, which cannot be
normalized.

\subsection{Why the characteristic equation turns calculus into algebra}
\index{differential operator!characteristic polynomial}

Let \(\mathrm{D}=\mathrm{d}/\mathrm{d}x\). A constant-coefficient linear
differential equation can be written as a polynomial in the differentiation
operator:
\begin{equation}
  P(\mathrm{D})\psi=0,
\end{equation}
where, for example,
\begin{equation}
  P(\mathrm{D})=\mathrm{D}^2+k^2
\end{equation}
for the box interior. Exponentials are special because differentiation acts on
them by scalar multiplication:
\begin{equation}
  \mathrm{D}^j e^{rx}=r^j e^{rx}.
\end{equation}
Consequently,
\begin{equation}
  P(\mathrm{D})e^{rx}=P(r)e^{rx}.
\end{equation}
The trial function converts the differential operator \(\mathrm{D}\) into the
algebraic number \(r\), and converts the operator polynomial
\(P(\mathrm{D})\) into the ordinary polynomial \(P(r)\). The differential
equation is satisfied when
\begin{equation}
  P(r)=0.
\end{equation}
We have not discarded the calculus. We have chosen functions on which the
derivative has the simple eigenvalue action
\begin{equation}
  \mathrm{D}e^{rx}=re^{rx}.
\end{equation}
That eigenfunction property is the bridge from the differential equation to
the algebraic characteristic equation.

This construction parallels the matrix characteristic equation from
Chapter~\ref{ch:training-hermitian-eigensystems}. For a matrix, one seeks
numbers \(\lambda\) for which \(H-\lambda I\) loses invertibility. For a
constant-coefficient differential equation, one seeks numbers \(r\) for which
\(P(r)=0\), making \(e^{rx}\) a solution of \(P(\mathrm{D})\psi=0\). Both methods turn
an operator question into an algebraic root problem, although the operators and
root variables have different meanings.

The method has a boundary. If differential-equation coefficients depend on
\(x\), then substituting \(e^{rx}\) generally does not reduce the complete
problem to one polynomial in \(r\). Constant coefficients are what make the
simple characteristic equation possible here.

\subsection{Normalization and the eigenfunction basis}

Before normalization, the \(n\)th eigenfunction is
\begin{equation}
  \psi_n(x)
  =A_n\sin\!\left(\frac{n\pi(x-a)}{L}\right).
\end{equation}
Using
\begin{equation}
  \int_a^b
  \sin^2\!\left(\frac{n\pi(x-a)}{L}\right)\,\mathrm{d}x
  =\frac{L}{2},
\end{equation}
the normalization condition gives \(|A_n|=\sqrt{2/L}\). Choosing a real
positive phase convention,
\begin{equation}
  \psi_n(x)
  =\sqrt{\frac{2}{L}}
   \sin\!\left(\frac{n\pi(x-a)}{L}\right).
\end{equation}
The eigenfunctions obey
\begin{equation}
  \int_a^b\psi_m(x)^*\psi_n(x)\,\mathrm{d}x=\delta_{mn}.
\end{equation}
They form an orthonormal basis of \(L^2(a,b)\). Consequently, a general state
can be expanded as
\begin{equation}
  \psi(x)=\sum_{n=1}^{\infty}c_n\psi_n(x),
  \qquad
  c_n=\int_a^b\psi_n(x)^*\psi(x)\,\mathrm{d}x.
\end{equation}
This is the infinite-dimensional analogue of assembling matrix eigenvectors as
basis columns in Chapter~\ref{ch:training-hermitian-eigensystems}. Applying the
Hamiltonian scales each basis component by its energy:
\begin{equation}
  \hat H_{\mathrm{box}}\psi
  =\sum_{n=1}^{\infty}c_nE_n\psi_n.
\end{equation}
Each one-dimensional energy level is nondegenerate: apart from normalization
and phase, the two boundary conditions leave one eigenfunction for each
positive integer \(n\).

\subsection{Born's interpretation of the eigenfunction}
\index{Born rule|textbf}
\index{probability density}
\index{measurement!position}
\index{measurement!energy}

The eigenfunction \(\psi_n(x)\) is a probability amplitude, not a probability
by itself. Born's interpretation assigns the position probability density
\begin{equation}
  \rho_n(x)
  =|\psi_n(x)|^2
  =\frac{2}{L}
   \sin^2\!\left(\frac{n\pi(x-a)}{L}\right).
\end{equation}
For any subinterval \([c,d]\subseteq[a,b]\), the probability that an ideal
position measurement returns a value in that subinterval is
\begin{equation}
  \Pr\!\left(x\in[c,d]\mid\psi_n\right)
  =\int_c^d|\psi_n(x)|^2\,\mathrm{d}x.
\end{equation}
Normalization ensures
\begin{equation}
  \Pr\!\left(x\in[a,b]\mid\psi_n\right)
  =\int_a^b|\psi_n(x)|^2\,\mathrm{d}x
  =1.
\end{equation}
The wavefunction has physical dimension \(\text{length}^{-1/2}\) in one
spatial dimension, so \(|\psi_n|^2\) has dimension
\(\text{length}^{-1}\), as a position density must.

The density vanishes at the walls and at every interior node of the sine
function. For the \(n\)th eigenfunction, the nodes occur at
\begin{equation}
  x_j=a+\frac{jL}{n},
  \qquad
  j=0,1,\ldots,n.
\end{equation}
There are \(n-1\) interior nodes. A maximum of \(|\psi_n(x)|^2\) identifies a
position region with larger measurement probability; it is not a classical
trajectory or a claim that the particle follows the plotted curve.

Multiplying an eigenfunction by a constant phase does not alter any position
probability:
\begin{equation}
  \left|e^{i\alpha}\psi_n(x)\right|^2
  =|\psi_n(x)|^2.
\end{equation}
This makes the phase freedom of an individual eigenvector physically visible:
the coordinate representation changes, while the position distribution does
not.

Born's rule also applies to the energy basis. For a normalized general state
\begin{equation}
  \psi(x)=\sum_{n=1}^{\infty}c_n\psi_n(x),
  \qquad
  \sum_{n=1}^{\infty}|c_n|^2=1,
\end{equation}
an ideal energy measurement returns \(E_n\) with probability
\begin{equation}
  \Pr(E_n\mid\psi)
  =|c_n|^2
  =\left|\langle\psi_n\vert\psi\rangle\right|^2.
\end{equation}
A state may therefore have a continuous position distribution and, at the same
time, a discrete distribution over energy outcomes. The chosen observable
determines which eigenbasis supplies the alternatives.

For a single stationary energy eigenstate, time evolution contributes only a
phase:
\begin{equation}
  \psi_n(x,t)
  =e^{-iE_nt/\hbar}\psi_n(x).
\end{equation}
Its position density is time independent because
\begin{equation}
  |\psi_n(x,t)|^2=|\psi_n(x)|^2.
\end{equation}
A superposition generally contains relative phases
\(e^{-i(E_n-E_m)t/\hbar}\), so its position density can change with time. A
full development of measurement and time evolution requires the broader
postulates of quantum mechanics; the statements here provide the interpretation
needed for the box eigensystem~\cite{shankar1994}.

\subsection{The Dirichlet Laplacian on the closed interval \([a,b]\)}
\index{Laplacian!one-dimensional Dirichlet}
\index{Dirichlet Laplacian}

In one Cartesian dimension, the scalar Laplacian is the second derivative:
\begin{equation}
  \Delta=\frac{\mathrm{d}^2}{\mathrm{d}x^2}.
\end{equation}
To specify an operator on the interval, the differential expression must be
combined with boundary conditions. The Dirichlet Laplacian uses
\begin{equation}
  \psi(a)=\psi(b)=0.
\end{equation}
The positive operator used in the kinetic energy is the negative Dirichlet
Laplacian,
\begin{equation}
  \hat L_D=-\Delta_D=-\frac{\mathrm{d}^2}{\mathrm{d}x^2}.
\end{equation}
In formal domain notation,
\begin{equation}
  \mathcal D(\hat L_D)
  =H^2(a,b)\cap H_0^1(a,b).
\end{equation}
Here the Sobolev-space notation records square-integrable derivatives and zero
boundary traces. The essential point for the present construction is that the
endpoints and Dirichlet conditions belong to the operator definition.

The continuum box Hamiltonian factors into a physical coefficient and this
geometric operator:
\begin{equation}
  \hat H_{\mathrm{box}}
  =\frac{\hbar^2}{2m}\hat L_D.
\end{equation}
The eigenvalues of \(\hat L_D\) are
\begin{equation}
  \mu_n=\left(\frac{n\pi}{L}\right)^2,
\end{equation}
and the physical energies are
\begin{equation}
  E_n=\frac{\hbar^2}{2m}\mu_n.
\end{equation}
This separates the interval geometry and boundary condition from the mass and
Planck-constant scaling.

\subsection{Discretizing the Dirichlet Laplacian as an operator}
\index{finite difference!particle in a box}
\index{discrete Laplacian!Dirichlet interval}

Choose \(N\) interior grid points on the closed interval \([a,b]\):
\begin{equation}
  x_j=a+jh,
  \qquad
  h=\frac{b-a}{N+1}=\frac{L}{N+1},
  \qquad
  j=1,\ldots,N.
\end{equation}
The endpoints are
\begin{equation}
  x_0=a,
  \qquad
  x_{N+1}=b.
\end{equation}
Dirichlet data fix
\begin{equation}
  \psi(x_0)=\psi(x_{N+1})=0,
\end{equation}
so those two values are not independent components of the numerical state
vector. Define the interior sample vector
\begin{equation}
  \boldsymbol\psi_h
  =
  \begin{bmatrix}
    \psi(x_1) & \psi(x_2) & \cdots & \psi(x_N)
  \end{bmatrix}^{\mathsf T}
  \in\mathbb C^N.
\end{equation}

At an interior point, the centered second-difference formula is
\begin{equation}
  \psi''(x_j)
  \approx
  \frac{\psi(x_{j+1})-2\psi(x_j)+\psi(x_{j-1})}{h^2}.
\end{equation}
Therefore the discrete second-derivative operator is
\begin{equation}
  D_{2,h}
  =\frac{1}{h^2}
  \begin{bmatrix}
    -2 & 1 & 0 & \cdots & 0\\
    1 & -2 & 1 & \ddots & \vdots\\
    0 & 1 & -2 & \ddots & 0\\
    \vdots & \ddots & \ddots & \ddots & 1\\
    0 & \cdots & 0 & 1 & -2
  \end{bmatrix}.
\end{equation}
The finite representation of the positive operator \(-\Delta_D\) is
\begin{equation}
  L_{D,h}=-D_{2,h}
  =\frac{1}{h^2}
  \begin{bmatrix}
    2 & -1 & 0 & \cdots & 0\\
    -1 & 2 & -1 & \ddots & \vdots\\
    0 & -1 & 2 & \ddots & 0\\
    \vdots & \ddots & \ddots & \ddots & -1\\
    0 & \cdots & 0 & -1 & 2
  \end{bmatrix}
  \in\mathbb R^{N\times N}.
\end{equation}
Finally, the discretized box Hamiltonian is
\begin{equation}
  H_h=\frac{\hbar^2}{2m}L_{D,h}.
\end{equation}
The three symbols \(D_{2,h}\), \(L_{D,h}\), and \(H_h\) represent different
objects: the second derivative, the positive operator representing the
negative Laplacian, and the kinetic Hamiltonian. Their signs and units must not
be collapsed.

\subsection{Eigenpairs of the discretized operator}

The discrete Dirichlet eigenvectors have components
\begin{equation}
  v_n(j)
  =\sqrt{\frac{2}{N+1}}
   \sin\!\left(\frac{jn\pi}{N+1}\right),
  \qquad
  j,n=1,\ldots,N.
\end{equation}
To see why, set
\begin{equation}
  \theta_n=\frac{n\pi}{N+1}
\end{equation}
and temporarily ignore normalization. The \(j\)th component of
\(L_{D,h}\mathbf v_n\) is
\begin{align}
  (L_{D,h}\mathbf v_n)_j
  &=\frac{1}{h^2}
    \left(
      2\sin(j\theta_n)
      -\sin((j-1)\theta_n)
      -\sin((j+1)\theta_n)
    \right)\\
  &=\frac{2-2\cos\theta_n}{h^2}\sin(j\theta_n)\\
  &=\frac{4}{h^2}\sin^2\!\left(\frac{\theta_n}{2}\right)
    \sin(j\theta_n).
\end{align}
The endpoint conditions are built into the sine sequence:
\begin{equation}
  v_n(0)=0,
  \qquad
  v_n(N+1)=\sin(n\pi)=0.
\end{equation}
Hence
\begin{equation}
  L_{D,h}\mathbf v_n
  =\mu_{n,h}\mathbf v_n,
\end{equation}
with discrete Laplacian eigenvalues
\begin{equation}
  \mu_{n,h}
  =\frac{4}{h^2}
   \sin^2\!\left(\frac{n\pi}{2(N+1)}\right).
\end{equation}
The Hamiltonian eigenvalues are therefore
\begin{equation}
  \varepsilon_n
  =\frac{\hbar^2}{2m}\mu_{n,h}
  =\frac{2\hbar^2}{mh^2}
   \sin^2\!\left(\frac{n\pi}{2(N+1)}\right).
\end{equation}

The \(N\) discrete eigenvectors form an orthonormal basis of \(\mathbb R^N\).
The infinitely many continuum sine functions form an orthonormal basis of
\(L^2(a,b)\). For a fixed mode \(n\), their normalization conventions are
related by the quadrature factor
\begin{equation}
  v_n(j)
  =\sqrt{h}\,\psi_n(x_j).
\end{equation}
For each fixed \(n\), \(\mu_{n,h}\to\mu_n\) and
\(\varepsilon_n\to E_n\) as \(h\to0\). This statement must not be strengthened
into uniform convergence of every resolvable mode as \(N\) grows;
Appendix~\ref{app:particle-box} examines that distinction with retained
calculated evidence.

\section{Finite Hamiltonian representations}

Let \(\mathcal S\) be a finite state space with ordered orthonormal basis
\(\mathcal B=\{\lvert b_i\rangle\}_{i=0}^{N-1}\). Once a Hamiltonian acting on
\(\mathcal S\) has been defined, its represented matrix is
\begin{equation}
  H^{(\mathcal B)}_{ij}
  =\langle b_i\lvert\hat H_{\mathcal S}\rvert b_j\rangle.
\end{equation}
If the decomposition has been transported consistently to the same space and
basis, then
\begin{equation}
  H^{(\mathcal B)}=T^{(\mathcal B)}+V^{(\mathcal B)}.
\end{equation}
Without the common-space and common-basis qualification, this array equation
has no justified operator interpretation.

\begin{pythonexample}{diagonalizing the one-dimensional box matrix}
import numpy as np

# Use dimensionless pedagogical units and declare both interval endpoints.
point_count = 8
left_endpoint = -0.5
right_endpoint = 1.5
length = right_endpoint - left_endpoint
mass = 1.0
hbar = 1.0
spacing = length / (point_count + 1)

# Assemble the centered Dirichlet Hamiltonian on interior grid points.
diagonal = 2.0 * np.ones(point_count, dtype=np.float64)
off_diagonal = -np.ones(point_count - 1, dtype=np.float64)
kinetic_core = (
    np.diag(diagonal)
    + np.diag(off_diagonal, k=1)
    + np.diag(off_diagonal, k=-1)
)
H_box = hbar**2 * kinetic_core / (2.0 * mass * spacing**2)

# eigh returns ascending energies and normalized eigenvectors as columns.
energies, states = np.linalg.eigh(H_box)
mode_indices = np.arange(1, point_count + 1, dtype=np.float64)
expected_discrete = (
    2.0
    * hbar**2
    / (mass * spacing**2)
    * np.sin(mode_indices * np.pi / (2.0 * (point_count + 1))) ** 2
)
np.testing.assert_allclose(energies, expected_discrete)

# Sample the shifted continuum sine basis with its quadrature factor.
grid = left_endpoint + spacing * mode_indices
sampled_sines = np.sqrt(spacing) * np.sqrt(2.0 / length) * np.sin(
    np.outer(grid - left_endpoint, mode_indices * np.pi / length)
)

# Eigenvector signs are arbitrary, so compare absolute basis overlaps.
overlaps = np.abs(states.conjugate().T @ sampled_sines)
np.testing.assert_allclose(overlaps, np.eye(point_count), atol=1.0e-12)
\end{pythonexample}

The arrays are synthetic teaching data in declared dimensionless units. The
checks establish exact identities of this finite matrix and sampled sine basis;
they do not establish a material energy or a general convergence result.

\documentcallout{Notebook to do}{Create
\path{examples/tutorials/research-monograph/foundations/quantum_operators.ipynb}
with explicit imports, explanatory inline comments, normalized states,
expectation values, the one-dimensional particle-in-a-box derivation and finite
matrix, and a clear separation between the continuum operator and its finite
representation. Keep it unexecuted in the repository.}

\section{What changes in an approximation}

Several distinct operations can produce a smaller or simpler Hamiltonian:
selecting a different physical model, projecting onto a retained subspace,
changing basis, discretizing a differential expression, truncating matrix
entries, or fitting parameters in an admissible model class. A correct account
names the operation rather than calling every output an ``effective
Hamiltonian.''

Chapter~\ref{ch:training-particle-box-higher-dimensions} next extends the box
to two and three dimensions and explains how Kronecker products represent
separate coordinate actions. Chapter~\ref{ch:training-periodic-geometry} then
supplies the geometry needed by periodic quantum representations, before later
chapters introduce Bloch-periodic finite differences, retained spaces,
isolated-band Fourier reduction, and evidence vocabulary.

\section{Exercises}

\begin{enumerate}
  \item Explain why an operator domain is not recoverable from the entries of a
        finite matrix alone.
  \item Show that the expectation value is unchanged under a consistent unitary
        basis transformation of both the state coordinates and operator matrix.
  \item Starting from the stationary Schrödinger equation, derive the permitted
        wave numbers and energies of the one-dimensional box.
  \item Normalize the shifted sine eigenfunctions and explain why they form a
        basis of \(L^2(a,b)\).
  \item Use Born's rule to compute the probability of finding the particle in a
        chosen subinterval \([c,d]\subseteq[a,b]\), and distinguish this
        position probability from the probability of an energy outcome.
  \item Derive \(D_{2,h}\), \(L_{D,h}\), and \(H_h\), explaining their signs,
        units, and represented meanings.
  \item Derive the discrete sine eigenvectors and distinguish their finite basis
        of \(\mathbb R^N\) from the infinite continuum basis.
  \item Give an example in which \(V(\mathbf r)\) is local in position space but
        has off-diagonal entries in another basis.
  \item State the compatibility requirements needed before adding separately
        supplied finite kinetic and potential matrices.
  \item Distinguish a list of eigenvalues, an eigensystem, a represented
        Hamiltonian, and an abstract Hamiltonian.
\end{enumerate}

\section{Evidence boundary}

A Hermiticity check, normalized eigenvectors, and a small residual can verify
bounded properties of a finite calculation. They do not establish that the
state space, domain, boundary conditions, Hamiltonian terms, or physical model
are adequate. Those choices must be declared before stronger evidence can be
interpreted.
